\section{Track C: Nested expectations: models and estimators, Part I}\newpage

\begin{talk}
  {An Adaptive Sampling Algorithm for Level-set Approximation}% [1] talk title
  {Abdul-Latefe Haji-Ali}% [2] speaker name
  {Heriot-Watt University}% [3] affiliations
  {a.hajiali@hw.ac.uk}% [4] email
  {Matteo Croci and Ian CJ Powell}% [5] coauthors
  {}% [6] special session. Leave this field empty for contributed talks.
    

  Let $D\subset\joinrel\subset\rm{I\!R}$$^d$ be a $d$-dimensional domain with compact
closure. We consider the approximation of the $d-1$ dimensional zero level-set
$\mathcal{L}_0 := \{x \in \overline{D} : f(x) = 0\}$ where the Lipschitz function $f$ is either
accessible directly or when $f(x) = \mathbb{E}$$\left[\tilde{f}(x)\right]$ for
all $x \in \overline{D}$. Given an approximation scheme with a priori error bounds
and $L^p$ bounds on the $\tilde{f}$-approximation error, we propose a
grid-based adaptive sampling scheme which produces an approximation to $\mathcal{L}_0$
with expected cost-complexity reduction of $\varepsilon^{\left(\frac{p+1}{\alpha p}\right)}$
compared to a non-adaptive scheme, where $\alpha$ is the known convergence rate of
an interpolation scheme. We provide the numerical analysis and experiments to
show that these rates hold in practice.

\medskip

% If you would like to include references, please do so by creating a simple list numbered by [1], [2], [3], \ldots. See example below.
% Please do not use the \texttt{bibliography} environment or \texttt{bibtex} files.
% APA reference style is recommended.
% \begin{enumerate}
%  \item[{[1]}] Niederreiter, Harald (1992). {\it Random number generation and quasi-Monte Carlo methods}. Society for Industrial and Applied Mathematics (SIAM).
%  \item[{[2]}] Roberts, Gareth O, \& Rosenthal, Jeffrey S. (2002).  Optimal scaling for various Metropolis-Hastings algorithms, \textbf{16}(4), 351--367.
% \end{enumerate}

% Equations may be used if they are referenced. Please note that the equation numbers may be different (but will be cross-referenced correctly) in the final program book.
\end{talk}

\begin{talk}
  {Posterior-Free A-Optimal Bayesian Design of Experiments via Conditional Expectation}% [1] talk title
  {Vinh Hoang}% [2] speaker name
  {RWTH-Aachen University}% [3] affiliations
  {hoang@uq.rwth-aachen.de}% [4] email
  {Luis Espath, Sebastian Krumscheid, Ra\'ul Tempone}% [5] coauthors
  
    
   

\medskip

We propose a novel approach for solving the A-optimal Bayesian design of experiments that does not require sampling or approximating the posterior distribution. In this setting, the objective function is the expected conditional variance (ECV).
Our method estimates the ECV by leveraging conditional expectation, which we approximate using its orthogonal projection property. We derive an asymptotic error bound for this estimator and validate it through numerical experiments.
The method is particularly efficient when the design parameter space is continuous. In such scenarios, the conditional expectation can be approximated non-locally using tools such as neural networks. To reduce the number of evaluations of the measurement model, we incorporate transfer learning and data augmentation.
Numerical results show that our method significantly reduces model evaluations compared to standard importance sampling-based techniques.
Code available at: \href{https://github.com/vinh-tr-hoang/DOEviaPACE}{https://github.com/vinh-tr-hoang/DOEviaPACE}.

\begin{enumerate}
    \item [{[1]}] Hoang, V., Espath, L., Krumscheid, S., \& Tempone, R. (2025).  
    Scalable method for Bayesian experimental design without integrating over posterior distribution. {\it SIAM ASA Journal on Uncertainty Quantification, 13}(1), 114-139. \\
    \href{https://doi.org/10.1137/23M1603364}{https://doi.org/10.1137/23M1603364}
\end{enumerate}


\end{talk}

\begin{talk}
  {QMC for Bayesian optimal experimental design with application to inverse problems governed by PDEs}% [1] talk title
  {Vesa Kaarnioja}% [2] speaker name
  {Free University of Berlin}% [3] affiliations
  {vesa.kaarnioja@fu-berlin.de}% [4] email
  {Claudia Schillings}% [5] coauthors
  {Nested expectations: models and estimators, Part I}% [6] special session. Leave this field empty for contributed talks. 
    
   
The goal in Bayesian optimal experimental design (OED) is to maximize the expected information gain for the reconstruction of unknown quantities in an experiment by optimizing the placement of measurements. The objective function in the resulting optimization problem involves a multivariate double integral over the high-dimensional parameter and data domains. For the efficient approximation of these integrals, we consider a sparse tensor product combination of quasi-Monte Carlo (QMC) cubature rules over the parameter and data domains. For the parameterization of the unknown quantities, we consider a model recently studied by Chernov and L\^{e} [1,2] as well as Harbrecht, Schmidlin, and Schwab [3] in which the input random field is assumed to belong to a Gevrey class. The Gevrey class contains functions that are infinitely many times continuously differentiable with a growth condition on the higher-order partial derivatives, but which are not analytic in general. Using the techniques developed in [4], we investigate efficient Bayesian OED for inverse problems governed by partial differential equations (PDEs).
\begin{enumerate}
 \item[{[1]}] Chernov, Alexey, \& L\^{e}, T\`{u}ng (2024). Analytic and Gevrey class regularity for parametric elliptic eigenvalue problems and applications. \emph{SIAM Journal on Numerical Analysis}, \textbf{62}(4), 1874--1900.
 \item[{[2]}] Chernov, Alexey, \& L\^{e}, T\`{u}ng (2024). Analytic and Gevrey class regularity for parametric semilinear reaction-diffusion problems and applications in uncertainty quantification. \emph{Computers \& Mathematics with Applications}, \textbf{164}, 116--130.
 \item[{[3]}] Harbrecht, Helmut, Schmidlin, Marc, \& Schwab, Christoph (2024). The Gevrey class implicit mapping theorem with applications to UQ of semilinear elliptic PDEs. \emph{Mathematical Models and Methods in Applied Sciences}, \textbf{34}(5), 881--917.
 \item[{[4]}] Kaarnioja, Vesa, \& Schillings, Claudia (2024). Quasi-Monte Carlo for Bayesian design of experiment problems governed by parametric PDEs. Preprint, \emph{arXiv:2405.03529 [math.NA]}.
\end{enumerate}

\end{talk}
